% Generated from locale/tr/content/proof-theory/normalization/permutations.tex; do not hand-edit.


\olfileid[tr]{pt}{nor}{per}

\olsection{Permütasyon Dönüştürmeleri}

\Log{N1i} sistemindeki kesmeler, !!a{elimination} kuralının büyük
öncülünde biten kesitlerdir. Normalleştirme ispatındaki durumlardan
biri, maksimal kesmelerin uzunluğunu ve böylece !!{proof} için kesme
uzunluğunu azaltmaktır. Uzunluğu $>1$ olan bir kesme, kesmedeki son
$!A$ !!{formula} oluşumunun bir çıkarımın (\Elim{\lor} ya da
\Elim{\lexists}) sonucu olup olmamasına ve hangi !!{elimination}
çıkarımının büyük öncülü olduğuna bağlı olarak birçok biçimde bitebilir.

Örneğin, ilgili çıkarımlar \Elim{\lor} ve \Elim{\land}, $!A \ident !B
\land !C$ ise ve aşağıdaki gibi biten bir $\delta_1$ alt-!!{proof}
varsa:
\[
\AxiomC{}
\RightLabel{$\delta_2$}
\DeduceC{$!D \lor !E$}
\AxiomC{$\Discharge{!D}{x}$}
\RightLabel{$\delta_3$}
\DeduceC{$!B \land !C$}
\AxiomC{$\Discharge{!E}{y}$}
\RightLabel{$\delta_4$}
\DeduceC{$!B \land !C$}
\DischargeRule{\Elim{\lor}}{x\,y}
\TrinaryInfC{$!B \land !C$}
\RightLabel{\Elim{\land}[1]}
\UnaryInfC{$!B$}
\DisplayProof
\]
Kesmedeki son $!A_n$ !!{formula} oluşumu, \Elim{\lor} çıkarımının sonucu
olan alttaki $!B \land !C$'dir. Sondan bir önceki $!A_{n-1}$
!!{formula} oluşumu ise \Elim{\lor} çıkarımının küçük öncüllerinden
biridir.

Bu kesmenin uzunluğunu azaltmak için \Elim{\lor} ve \Elim{\land}
çıkarımlarının sırasını değiştirebilir (``permüte edebiliriz''). Başka
bir deyişle, $\delta$ içindeki $\delta_1$ alt ispatını aşağıdakiyle
değiştiririz:
\[
\AxiomC{}
\RightLabel{$\delta_2$}
\DeduceC{$!D \lor !E$}
\AxiomC{$\Discharge{!D}{x}$}
\RightLabel{$\delta_3$}
\DeduceC{$!B \land !C$}
\RightLabel{\Elim{\land}[1]}
\UnaryInfC{$!B$}
\AxiomC{$\Discharge{!E}{y}$}
\RightLabel{$\delta_4$}
\DeduceC{$!B \land !C$}
\RightLabel{\Elim{\land}[1]}
\UnaryInfC{$!B$}
\DischargeRule{\Elim{\lor}}{x\,y}
\TrinaryInfC{$!B$}
\DisplayProof
\]
Önceden \Elim{\lor} altındaki \Elim{\land} öncülünde biten kesmeler,
şimdi \Elim{\lor} çıkarımının küçük öncülleri üzerindeki \Elim{\land}
çıkarımlarından birinin öncülünde biter; yani her birinin uzunluğu
$1$ azalmıştır.

Ayrıca, bütünüyle $\delta_2$, $\delta_3$ ya da $\delta_4$ içinde veya
$\delta$ içindeki $\delta_1$'in bütünüyle dışında kalan her kesme
değişmeden kalır: Yine aynı !!{formula} biçimlerini içerir ve aynı çıkarımlardan
geçer; dolayısıyla özellikle, $\delta$ içindeki karşılık gelen kesmeyle
aynı dereceye ve uzunluğa sahiptir. Böylece yeni !!{proof} için kesme
derecesi artmamış, yeni !!{proof} için kesme uzunluğu ise eski
!!{proof} için kesme
uzunluğundan daha küçük olmuştur.

\begin{prob}
Bir permütasyon dönüştürmesinden sonra en az bir kesmenin uzunluğu
$1$ azalır. \Elim{\land} çıkarımını bir \Elim{\lor} üzerinden permüte
edersek aslında uzunluğu azalan iki kesme kesiti vardır; dolayısıyla bu
durumda !!{proof} için kesme uzunluğu en az $2$ azalır. !!a{proof} için
kesme uzunluğu bu tür tek bir permütasyonla $2$'den fazla azalabilir
mi? Kesme \emph{derecesi} azalabilir mi?
\end{prob}

!!a{elimination} kuralını bir \Elim{\lor} ya da \Elim{\lexists}
kuralının üzerine taşıma işlemine \emph{permütasyon dönüştürmesi}
denir. Bazı kural çiftlerine ilişkin kalıplar
\olref{tab:permutation-conversions}'da verilmiştir. (Kalanlar alıştırma
olarak bırakılmıştır.)

% \begin{table}
% \begin{multline*}
% \shoveleft{\AxiomC{$!D \lor !E$}
% \AxiomC{$\Discharge{!D}{x}$}
% \shortDeduce
% \DeduceC{$!B \land !C$}
% \AxiomC{$\Discharge{!E}{y}$}
% \shortDeduce
% \DeduceC{$!B \land !C$}
% \DischargeRule{\Elim{\lor}}{x\,y}
% \TrinaryInfC{$!B \land !C$}
% \RightLabel{\Elim{\land}[1]}
% \UnaryInfC{$!B$}
% \DisplayProof}\\
% \shoveright{\longrightarrow\ 
% \AxiomC{$!D \lor !E$}
% \AxiomC{$\Discharge{!D}{x}$}
% \shortDeduce
% \DeduceC{$!B \land !C$}
% \RightLabel{\Elim{\land}[1]}
% \UnaryInfC{$!B$}
% \AxiomC{$\Discharge{!E}{y}$}
% \shortDeduce
% \DeduceC{$!B \land !C$}
% \RightLabel{\Elim{\land}[1]}
% \UnaryInfC{$!B$}
% \DischargeRule{\Elim{\lor}}{x\,y}
% \TrinaryInfC{$!B$}
% \DisplayProof}
% \\
% \shoveleft{\AxiomC{$!D \lor !E$}
% \AxiomC{$\Discharge{!D}{x}$}
% \shortDeduce
% \DeduceC{$!B \lif !C$}
% \AxiomC{$\Discharge{!E}{y}$}
% \shortDeduce
% \DeduceC{$!B \lif !C$}
% \DischargeRule{\Elim{\lor}}{x\,y}
% \TrinaryInfC{$!B \lif !C$}
% \AxiomC{$!B$}
% \RightLabel{\Elim{\lif}}
% \BinaryInfC{$!C$}
% \DisplayProof}\\
% \shoveright{\longrightarrow\ 
% \AxiomC{$!D \lor !E$}
% \AxiomC{$\Discharge{!D}{x}$}
% \shortDeduce
% \DeduceC{$!B \lif !C$}
% \AxiomC{$!B$}
% \RightLabel{\Elim{\lif}}
% \BinaryInfC{$!C$}
% \AxiomC{$\Discharge{!E}{y}$}
% \shortDeduce
% \DeduceC{$!B \lif !C$}
% \AxiomC{$!B$}
% \RightLabel{\Elim{\lif}}
% \BinaryInfC{$!C$}
% \DischargeRule{\Elim{\lor}}{x\,y}
% \TrinaryInfC{$!C$}
% \DisplayProof}
% \\
% \shoveleft{\AxiomC{$!D \lor !E$}
% \AxiomC{$\Discharge{!D}{x}$}
% \shortDeduce
% \DeduceC{$!B \lif !C$}
% \AxiomC{$\Discharge{!E}{y}$}
% \shortDeduce
% \DeduceC{$!B \lif !C$}
% \DischargeRule{\Elim{\lor}}{x\,y}
% \TrinaryInfC{$!B \lif !C$}
% \AxiomC{$!B$}
% \RightLabel{\Elim{\lif}}
% \BinaryInfC{$!C$}
% \DisplayProof}\\
% \shoveright{\longrightarrow\ 
% \AxiomC{$!D \lor !E$}
% \AxiomC{$\Discharge{!D}{x}$}
% \shortDeduce
% \DeduceC{$!B \lif !C$}
% \AxiomC{$!B$}
% \RightLabel{\Elim{\lif}}
% \BinaryInfC{$!C$}
% \AxiomC{$\Discharge{!E}{y}$}
% \shortDeduce
% \DeduceC{$!B \lif !C$}
% \AxiomC{$!B$}
% \RightLabel{\Elim{\lif}}
% \BinaryInfC{$!C$}
% \DischargeRule{\Elim{\lor}}{x\,y}
% \TrinaryInfC{$!C$}
% \DisplayProof}
% \\
% \shoveleft{\AxiomC{$!D \lor !E$}
% \AxiomC{$\Discharge{!D}{x}$}
% \shortDeduce
% \DeduceC{$\lforall[x][!B(x)]$}
% \AxiomC{$\Discharge{!E}{y}$}
% \shortDeduce
% \DeduceC{$\lforall[x][!B(x)]$}
% \DischargeRule{\Elim{\lor}}{x\,y}
% \TrinaryInfC{$\lforall[x][!B(x)]$}
% \RightLabel{\Elim{\lforall}}
% \UnaryInfC{$!B(t)$}
% \DisplayProof}\\
% \shoveright{\longrightarrow\ 
% \AxiomC{$!D \lor !E$}
% \AxiomC{$\Discharge{!D}{x}$}
% \shortDeduce
% \DeduceC{$\lforall[x][!B(x)]$}
% \RightLabel{\Elim{\lforall}}
% \UnaryInfC{$!B(t)$}
% \AxiomC{$\Discharge{!E}{y}$}
% \shortDeduce
% \DeduceC{$\lforall[x][!B(x)]$}
% \RightLabel{\Elim{\lforall}}
% \UnaryInfC{$!B(t)$}
% \DischargeRule{\Elim{\lor}}{x\,y}
% \TrinaryInfC{$!B(t)$}
% \DisplayProof}
% \end{multline*}
% \caption{Some permutation conversions for \Log{N1i}.}
% \ollabel{tab:permutation-conversions}
% \end{table}

{\small
\begin{longtable}{@{}l@{}}
\AxiomC{$!D \lor !E$}
\AxiomC{$\Discharge{!D}{x}$}
\shortDeduce
\DeduceC{$!B \land !C$}
\AxiomC{$\Discharge{!E}{y}$}
\shortDeduce
\DeduceC{$!B \land !C$}
\DischargeRule{\Elim{\lor}}{x\,y}
\TrinaryInfC{$!B \land !C$}
\RightLabel{\Elim{\land}[1]}
\UnaryInfC{$!B$}
\DisplayProof
$\longrightarrow$\\[6ex]
\quad\AxiomC{$!D \lor !E$}
\AxiomC{$\Discharge{!D}{x}$}
\shortDeduce
\DeduceC{$!B \land !C$}
\RightLabel{\Elim{\land}[1]}
\UnaryInfC{$!B$}
\AxiomC{$\Discharge{!E}{y}$}
\shortDeduce
\DeduceC{$!B \land !C$}
\RightLabel{\Elim{\land}[1]}
\UnaryInfC{$!B$}
\DischargeRule{\Elim{\lor}}{x\,y}
\TrinaryInfC{$!B$}
\DisplayProof
\\[10ex]
\AxiomC{$!D \lor !E$}
\AxiomC{$\Discharge{!D}{x}$}
\shortDeduce
\DeduceC{$!B \lif !C$}
\AxiomC{$\Discharge{!E}{y}$}
\shortDeduce
\DeduceC{$!B \lif !C$}
\DischargeRule{\Elim{\lor}}{x\,y}
\TrinaryInfC{$!B \lif !C$}
\AxiomC{$!B$}
\RightLabel{\Elim{\lif}}
\BinaryInfC{$!C$}
\DisplayProof
$\longrightarrow$\\[6ex]
\AxiomC{$!D \lor !E$}
\AxiomC{$\Discharge{!D}{x}$}
\shortDeduce
\DeduceC{$!B \lif !C$}
\AxiomC{$!B$}
\RightLabel{\Elim{\lif}}
\BinaryInfC{$!C$}
\AxiomC{$\Discharge{!E}{y}$}
\shortDeduce
\DeduceC{$!B \lif !C$}
\AxiomC{$!B$}
\RightLabel{\Elim{\lif}}
\BinaryInfC{$!C$}
\DischargeRule{\Elim{\lor}}{x\,y}
\TrinaryInfC{$!C$}
\DisplayProof
\\[10ex]
\AxiomC{$!D \lor !E$}
\AxiomC{$\Discharge{!D}{x}$}
\shortDeduce
\DeduceC{$\lforall[x][!B(x)]$}
\AxiomC{$\Discharge{!E}{y}$}
\shortDeduce
\DeduceC{$\lforall[x][!B(x)]$}
\DischargeRule{\Elim{\lor}}{x\,y}
\TrinaryInfC{$\lforall[x][!B(x)]$}
\RightLabel{\Elim{\lforall}}
\UnaryInfC{$!B(t)$}
\DisplayProof
\quad$\longrightarrow$\\[8ex]
\quad\AxiomC{$!D \lor !E$}
\AxiomC{$\Discharge{!D}{x}$}
\shortDeduce
\DeduceC{$\lforall[x][!B(x)]$}
\RightLabel{\Elim{\lforall}}
\UnaryInfC{$!B(t)$}
\AxiomC{$\Discharge{!E}{y}$}
\shortDeduce
\DeduceC{$\lforall[x][!B(x)]$}
\RightLabel{\Elim{\lforall}}
\UnaryInfC{$!B(t)$}
\DischargeRule{\Elim{\lor}}{x\,y}
\TrinaryInfC{$!B(t)$}
\DisplayProof
\\
\caption{\Log{N1i} için bazı permütasyon dönüştürmeleri.}
\ollabel{tab:permutation-conversions}
\end{longtable}
}

\begin{prob}
Ardından \Elim{\lor} ya da \Elim{\lnot} gelen \Elim{\lor} için
permütasyon dönüştürmelerini verin.
\end{prob}

\begin{prob}
Ardından \Elim{\lexists} gelen \Elim{\lexists} için permütasyon
dönüştürmelerini verin. İkinci durumda hiçbir özdeğişken koşulunun
neden ihlal edilmediğini açıklayın.
\end{prob}

Bir permütasyon dönüştürmesi uyguladığımızda herhangi bir özdeğişken
koşulunu ihlal etmediğimizden emin olmalıyız. Şunu ele alalım:
\[
\AxiomC{}
\RightLabel{$\delta_2$}
\DeduceC{$!D \lor !E$}
\AxiomC{$\Discharge{!D}{x}$}
\RightLabel{$\delta_3$}
\DeduceC{$\lexists[x][!B(x)]$}
\AxiomC{$\Discharge{!E}{y}$}
\RightLabel{$\delta_4$}
\DeduceC{$\lexists[x][!B(x)]$}
\DischargeRule{\Elim{\lor}}{x\,y}
\TrinaryInfC{$\lexists[x][!B(x)]$}
\AxiomC{$\Discharge{!B(c)}{z}$}
\RightLabel{$\delta_5$}
\DeduceC{$!C$}
\DischargeRule{\Elim{\lexists}}{z}
\BinaryInfC{$!C$}
\DisplayProof
\]
Bunu aşağıdakiyle değiştirirsek
\[
\AxiomC{}
\RightLabel{$\delta_2$}
\DeduceC{$!D \lor !E$}
\AxiomC{$\Discharge{!D}{x}$}
\RightLabel{$\delta_3$}
\DeduceC{$\lexists[x][!B(x)]$}
\AxiomC{$\Discharge{!B(c)}{z}$}
\RightLabel{$\delta_5$}
\DeduceC{$!C$}
\DischargeRule{\Elim{\lexists}}{z}
\BinaryInfC{$!C$}
\AxiomC{$\Discharge{!E}{y}$}
\RightLabel{$\delta_4$}
\DeduceC{$\lexists[x][!B(x)]$}
\AxiomC{$\Discharge{!B(c)}{z}$}
\RightLabel{$\delta_5$}
\DeduceC{$!C$}
\DischargeRule{\Elim{\lexists}}{z}
\BinaryInfC{$!C$}
\DischargeRule{\Elim{\lor}}{x\,y}
\TrinaryInfC{$!C$}
\DisplayProof
\]
özdeğişken $c$'nin, örneğin $!D$ içinde geçebileceğinden kaygılanabiliriz.
Nitekim $!D$ varsayımı özgün $\Elim{\lexists}$ çıkarımının üzerinde
boşa çıkarılır; dolayısıyla özgün !!{proof} içinde, \Elim{\lexists}
çıkarımının büyük öncülü üzerinde açık olan bir varsayımda geçmez ve
özdeğişken koşulu sağlanır. Ne var ki yeni !!{proof} içinde $!D$,
\Elim{\lexists} çıkarımlarından sonrasına kadar boşa çıkarılmaz; bu
durumda özdeğişken koşulu ihlal edilmiş olurdu. Ancak ele aldığımız
!!{proof} örneklerinin düzenli olduğunu varsaydığımızı hatırlayalım: Her özdeğişken yalnızca
bir çıkarımın özdeğişkenidir ve !!{proof} içinde yalnızca
\Elim{\lexists} çıkarımının küçük öncülü üzerinde geçer. Özellikle $c$,
$\delta_3$ ya da $\delta_4$ içinde geçemez.

Bu örnek başka bir şeyi daha gösterir: Yeni !!{proof} içinde
$\delta_5$'in iki kopyası olduğuna dikkat edin. Bu, $\delta_5$ bir
maksimal kesme içeriyorsa kesme uzunluğunu \emph{azaltmamış} olduğumuz
anlamına gelir! Bu nedenle bir permütasyon dönüştürmesini yalnızca
$\delta_5$ içinde maksimal kesme bulunmadığını bildiğimizde uygulamaya
dikkat etmeliyiz. Bunu her zaman \emph{en üstteki} maksimal kesmeyi
seçerek sağlarız; yani alt-!!{proof} (bu durumda $\delta_1$), alt
!!{proof} içindeki son çıkarımın üzerinde biten hiçbir maksimal kesme
içermeyecek biçimde bir maksimal kesme seçeriz. Bu, $\delta_5$ içinde
(hatta $\delta_2$, $\delta_3$ ya da $\delta_4$ içinde) başka maksimal
kesmeler bulunması olasılığını dışlar.

